% Generated draft: requires mathematical and source review.
\documentclass[11pt]{article}
\usepackage[margin=25mm]{geometry}
\usepackage{fontspec}
\setmainfont{Latin Modern Roman}
\usepackage{amsmath,amssymb,mathtools,unicode-math}
\setmathfont{Latin Modern Math}
\usepackage{enumitem,needspace,xurl,hyperref}
\setlength{\parindent}{0pt}
\setlength{\parskip}{5pt}
\setlength{\emergencystretch}{3em}
\newcommand{\N}{\mathbb{N}}
\newcommand{\Z}{\mathbb{Z}}
\newcommand{\Q}{\mathbb{Q}}
\newcommand{\R}{\mathbb{R}}
\newcommand{\C}{\mathbb{C}}
\newcommand{\F}{\mathbb{F}}
\newcommand{\A}{\mathbb{A}}
\newcommand{\PP}{\mathbb{P}}
\newcommand{\E}{\mathbb{E}}
\newcommand{\eps}{\varepsilon}
\newcommand{\dd}{\,\mathrm{d}}
\begin{document}
\section*{UnsolvedMath ID: 30005146}
\textbf{Source problem code:} catalogue entry
\subsection{Definitions and mathematical statement}
In the usual Riemannian interpretation, an Einstein metric on a smooth manifold \(M\) is a Riemannian metric \(g\) satisfying
\[
\operatorname{Ric}(g)=\lambda g
\]
for a real constant \(\lambda\), where \(\operatorname{Ric}(g)\) denotes the Ricci curvature tensor. The supplied record does not explicitly specify the category of manifolds or metrics.

The source asks the following stability question: if a manifold \(M\) admits a metric \(\bar g\) that is almost Einstein in a suitable sense, does there exist an Einstein metric \(g\) on the same manifold that is close to \(\bar g\)? Here both the almost-Einstein condition and the required closeness remain unspecified. The dimension, compactness, completeness, boundary conditions, normalization of the Einstein constant, geometric bounds, and topology or norm for comparing metrics are also unspecified. Consequently, the record does not determine a fully quantified implication with a prescribed error threshold and conclusion.

The original statement additionally contains a damaged formula labelled \((4)\). A tentative reading of its visible symbols is
\[
\text{for every }x\in M,\qquad
 e^{b\,d(x,M_{\mathrm{thick}})}
 \int_M e^{-(2-\delta)d(x,y)}
 \bigl|\operatorname{Ric}(\bar g)+2\bar g\bigr|_{\bar g}^{2}(y)
 \,\mathrm{dvol}_{\bar g}(y)
 \leq \varepsilon^{2}.
\tag{4}
\]
This is a provisional transcription, not an established hypothesis of the general question. Under this reading, \(d\) would denote the distance induced by \(\bar g\), \(\mathrm{dvol}_{\bar g}\) its volume measure, and \(|\cdot|_{\bar g}\) its tensor norm. The source does not define \(M_{\mathrm{thick}}\), the parameters \(b,\delta,\varepsilon\), their ranges or dependencies, or the other hypotheses accompanying \((4)\). The residual \(\operatorname{Ric}(\bar g)+2\bar g\) uses Einstein constant \(-2\), but the record does not establish that this normalization applies to the unrestricted question. No precise stability theorem or unrestricted affirmative assertion is supplied.
\subsection{Short English statement}
The question asks whether a metric that is sufficiently close to being Einstein guarantees an Einstein metric nearby on the same manifold. The source leaves the hypotheses and meaning of closeness unresolved, and an accompanying weighted curvature inequality is damaged and incomplete. The supplied dated curation labels the problem partially solved under strong geometric hypotheses; this draft does not independently verify that assessment.
\subsection{Sources}
\begin{itemize}
\item Local registry: supplied UnsolvedMath record, permanent ID \(30005146\), in \texttt{data/unsolvedmath/problems.json}; exact source problem code \texttt{catalogue entry}; title: \emph{Stability of Almost-Einstein Metrics}.
\item Supplied report citation: \emph{Geometrie (2023)}, Oberwolfach Reports, DOI 10.4171/owr/2022/28, pp.~1551--1601. \url{https://doi.org/10.4171/owr/2022/28}. The record gives proposed year \(2022\) and the tag \texttt{owr-year:2022}; the report citation's displayed year is retained as supplied.
\item Attribution in the original extract: \emph{Stability of Einstein metrics}, Frieder Jäckel, joint work with Ursula Hamenstädt. The separate \texttt{proposed\_by} field is empty, so this attribution does not establish who proposed the general question.
\item Supplied literature reference: Hamenstadt--Jackel, stability and effective hyperbolization. \url{https://arxiv.org/abs/2206.10438}.
\item Supplied literature reference: Jackel, effective higher-dimensional stability. \url{https://arxiv.org/abs/2502.16692}.
\item Supplied literature reference: OWR report. \url{https://ems.press/content/serial-article-files/46965}.
\item Literature assessment checked at \texttt{2026-08-22}: supplied dated curation, not independently verified. It labels the problem \texttt{partially\_solved} and reports stability under strong geometric hypotheses: almost-hyperbolic metrics yield nearby Einstein metrics in dimension three, and effective negatively curved results cover several higher-dimensional ranges. It describes the unrestricted almost-Einstein question as open. The background explicitly calls this a dated literature triage derived from the supplied OWR curation and requests independent verification of time-sensitive status and cited claims.
\item The registry labels the statement \texttt{corrected\_verified} and says that only clear formula delimiters, notation, grammar, and typography were repaired, with mathematical content and scope preserved. These are supplied editorial claims and do not resolve the missing hypotheses or constitute independent mathematical review.
\end{itemize}
\end{document}
